\documentclass[12pt]{article}

\usepackage{amsmath}
\usepackage{amssymb}
\usepackage{enumitem}
\usepackage{mathtools}

\title{Shortest Path Problem in Currency Trading}
\author{Algorithms Assignment}
\date{\today}

\begin{document}

\maketitle

\section*{Shortest Path Algorithms for Currency Trading}

\begin{enumerate}[label=\textbf{4.\arabic*.}]
\setcounter{enumi}{20}

\item Shortest path algorithms can be applied in currency trading. Let $c_1, c_2, \ldots, c_n$ be various currencies;
for instance, $c_1$ might be dollars, $c_2$ pounds, and $c_3$ lire. For any two currencies $c_i$ and
$c_j$, there is an exchange rate $r_{i,j}$; this means that you can purchase $r_{i,j}$ units of currency $c_j$ in
exchange for one unit of $c_i$. These exchange rates satisfy the condition that $r_{i,j} \cdot r_{j,i} < 1$, so that if
you start with a unit of currency $c_i$, change it into currency $c_j$ and then convert back to currency
$c_i$, you end up with less than one unit of currency $c_i$ (the difference is the cost of the transaction).

\begin{enumerate}[label=(\alph*)]
\item Give an efficient algorithm for the following problem: Given a set of exchange rates $r_{i,j}$,
and two currencies $s$ and $t$, find the most advantageous sequence of currency exchanges for
converting currency $s$ into currency $t$. Toward this goal, you should represent the currencies
and rates by a graph whose edge lengths are real numbers.
\end{enumerate}

The exchange rates are updated frequently, reflecting the demand and supply of the various
currencies. Occasionally the exchange rates satisfy the following property: there is a sequence of
currencies $c_{i_1}, c_{i_2}, \ldots, c_{i_k}$ such that $r_{i_1,i_2} \cdot r_{i_2,i_3} \cdot \ldots \cdot r_{i_{k-1},i_k} \cdot r_{i_k,i_1} > 1$. This means that by starting
with a unit of currency $c_{i_1}$ and then successively converting it to currencies $c_{i_2}, c_{i_3}, \ldots, c_{i_k}$, and
finally back to $c_{i_1}$, you would end up with more than one unit of currency $c_{i_1}$. Such anomalies
last only a fraction of a minute on the currency exchange, but they provide an opportunity for
risk-free profits.

\begin{enumerate}[label=(\alph*)]
\setcounter{enumii}{1}
\item Give an efficient algorithm for detecting the presence of such an anomaly. Use the graph
representation you found above.
\end{enumerate}

\end{enumerate}

\end{document}